% Part: sets-functions-relations
% Chapter: relations
% Section: equivalence-relations

\documentclass[../../../include/open-logic-section]{subfiles}

\begin{document}

\olfileid{sfr}{rel}{eqv}

\olsection{তুল্যতা সম্পর্ক}

একটি সেটের উপর অভিন্নতার সম্পর্ক প্রতিবিম্ব ধর্মবিশিষ্ট, প্রতিসম এবং পরিযায়ী। এই তিনটি ধর্মই আছে এমন সম্পর্ক~$R$ খুবই সাধারণ।

\begin{defn}[তুল্যতা সম্পর্ক]
প্রতিবিম্ব ধর্মবিশিষ্ট, প্রতিসম এবং পরিযায়ী একটি সম্পর্ক $R \subseteq A^2$-কে \emph{তুল্যতা সম্পর্ক} বলা হয়। $A$-এর !!^{element}s $x$ ও $y$-কে \emph{$R$-তুল্য} বলা হয় যদি~$Rxy$ হয়।
\end{defn}

তুল্যতা সম্পর্ক থেকে \emph{তুল্যতা শ্রেণি} ধারণাটি আসে। একটি তুল্যতা সম্পর্ক ক্ষেত্রটিকে আলাদা আলাদা বিভাজনখণ্ডে ভাগ করে। প্রতিটি খণ্ডের ভিতরে সব বস্তু পরস্পরের সঙ্গে সম্পর্কিত; আর ভিন্ন খণ্ডের কোনো বস্তু একে অপরের সঙ্গে সম্পর্কিত নয়। কখনো কখনো এই খণ্ডগুলি নিয়ে \emph{সরাসরি} আলোচনা করা সুবিধাজনক। সেই উদ্দেশ্যে আমরা একটি সংজ্ঞা দিচ্ছি:

\begin{defn}\ollabel{def:equivalenceclass}
ধরা যাক, $R \subseteq A^2$ একটি তুল্যতা সম্পর্ক। প্রত্যেক $x \in A$-এর জন্য $A$-তে $x$-এর \emph{তুল্যতা শ্রেণি} হল সেট $\equivrep{x}{R} = \Setabs{y \in A}{Rxy}$। $R$ অনুসারে $A$-এর \emph{ভাগসেট} হল $\equivclass{A}{R} = \Setabs{\equivrep{x}{R}}{x \in A}$, অর্থাৎ এই তুল্যতা শ্রেণিগুলির সেট।
\end{defn}

পরের ফলটি প্রমাণ করে যে তুল্যতা শ্রেণিগুলি সত্যিই $A$-এর বিভাজনখণ্ড; এইভাবে এটি তুল্যতা শ্রেণির সংজ্ঞাটির যথার্থতা প্রতিষ্ঠা করে:

\begin{prop}
যদি $R \subseteq A^2$ একটি তুল্যতা সম্পর্ক হয়, তবে $Rxy$ যদি এবং কেবল যদি $\equivrep{x}{R} = \equivrep{y}{R}$।
\end{prop}

\begin{proof}
বাঁ দিক থেকে ডান দিকের জন্য ধরা যাক $Rxy$, এবং ধরা যাক $z \in \equivrep{x}{R}$। তাহলে সংজ্ঞা অনুসারে $Rxz$। যেহেতু $R$ তুল্যতা সম্পর্ক, তাই $Ryz$। (বিস্তারিতভাবে: $Rxy$ এবং~$R$ প্রতিসম বলে $Ryx$; আবার $Rxz$ এবং~$R$ পরিযায়ী বলে~$Ryz$।) তাই $z \in \equivrep{y}{R}$। সাধারণীকরণ করলে পাই $\equivrep{x}{R} \subseteq \equivrep{y}{R}$। কিন্তু ঠিক একইভাবে $\equivrep{y}{R} \subseteq \equivrep{x}{R}$। সুতরাং সদস্যভিত্তিক সমতার নীতি অনুসারে $\equivrep{x}{R} = \equivrep{y}{R}$।

ডান দিক থেকে বাঁ দিকের জন্য ধরা যাক $\equivrep{x}{R} = \equivrep{y}{R}$। যেহেতু $R$ প্রতিবিম্ব ধর্মবিশিষ্ট, তাই $Ryy$; ফলে $y \in \equivrep{y}{R}$। অতএব $\equivrep{x}{R} = \equivrep{y}{R}$ অনুমান থেকে $y \in \equivrep{x}{R}$-ও হয়। তাই $Rxy$।
\end{proof}

\begin{ex}
তুল্যতা সম্পর্কের একটি সুন্দর উদাহরণ পাওয়া যায় মডুলার পাটিগণিতে। যেকোনো $a$, $b$, এবং $n \in \PosInt$-এর জন্য বলি $a \equiv_n b$ যদি এবং কেবল যদি $a$-কে~$n$ দিয়ে ভাগ করলে এবং $b$-কে~$n$ দিয়ে ভাগ করলে একই ভাগশেষ পাওয়া যায়। (কিছুটা বেশি প্রতীকী ভাষায়: $a \equiv_n b$ যদি এবং কেবল যদি কোনো $k \in \Int$-এর জন্য $a - b = kn$ হয়।) এখন, যেকোনো~$n$-এর জন্য $\equiv_n$ একটি তুল্যতা সম্পর্ক। আর $\equiv_n$ থেকে ঠিক $n$টি ভিন্ন তুল্যতা শ্রেণি পাওয়া যায়; অর্থাৎ $\equivclass{\Nat}{\equiv_n}$-এ $n$টি !!{element}s আছে। এগুলি হল: $n$ দিয়ে নিঃশেষে বিভাজ্য সংখ্যাগুলির সেট, অর্থাৎ $\equivrep{0}{\equiv_n}$; $n$ দিয়ে ভাগ করলে ভাগশেষ~$1$ হয় এমন সংখ্যাগুলির সেট, অর্থাৎ $\equivrep{1}{\equiv_n}$; \ldots; এবং $n$ দিয়ে ভাগ করলে ভাগশেষ~$n-1$ হয় এমন সংখ্যাগুলির সেট, অর্থাৎ~$\equivrep{n-1}{\equiv_n}$।
\end{ex}

\begin{prob}
দেখাও যে যেকোনো $n \in \PosInt$-এর জন্য $\equiv_n$ একটি তুল্যতা সম্পর্ক, এবং $\equivclass{\Nat}{\equiv_n}$-এর সদস্যসংখ্যা ঠিক $n$।
\end{prob}

\end{document}
